\introduction % Структурный элемент: ВВЕДЕНИЕ

Современные программные системы всё чаще разворачиваются как совокупность сервисов, средств наблюдаемости и вспомогательных workload, распределённых между несколькими средами и управляемыми разными командами. В таких условиях инфраструктура перестаёт быть статическим набором серверов и конфигурационных файлов. Она становится изменяемой системой, состояние которой должно постоянно соответствовать архитектурным решениям и эксплуатационным требованиям.

Актуальность темы выпускной квалификационной работы обусловлена тем, что ручное сопровождение инфраструктуры плохо масштабируется при росте числа хостов и сервисов. Даже если исходная конфигурация была подготовлена корректно, со временем между проектным и фактическим состоянием возникают расхождения. Они появляются из-за ручных изменений, временной недоступности источников данных, отличий между окружениями и ошибок конфигурации. Такие расхождения принято рассматривать как drift: ситуация, при которой фактическое состояние инфраструктуры отличается от ожидаемого.

Проблема drift особенно существенна для эксплуатационных команд, поскольку она приводит к снижению предсказуемости системы. Инженер может считать, что на управляемом хосте должен быть запущен определённый workload, однако фактически он может отсутствовать, быть остановлен или не наблюдаться средствами инвентаризации. Без автоматизированного механизма сравнения целевого и фактического состояния такие проблемы обнаруживаются поздно: через инциденты, ручной аудит или косвенные симптомы в мониторинге.

Одним из способов снижения этих рисков является переход от разрозненных ручных операций к декларативной модели управления. В такой модели команда описывает желаемое состояние системы, а программный контур управления получает фактическое состояние, сравнивает его с желаемым и инициирует действия по устранению расхождений. Близкая идея лежит в основе Kubernetes-контроллеров: они наблюдают состояние объектов и выполняют действия, необходимые для приближения текущего состояния к заданному \cite{kubernetes_controllers}. Однако Kubernetes-ориентированные решения не всегда подходят для управления host-level workload вне кластера или для ситуаций, где источником целевого состояния является архитектурная модель системы.

В данной работе рассматривается подход Architecture-as-Code как способ превратить архитектурное описание в машинно обрабатываемый источник desired state. В отличие от обычной документации, которая часто используется только для коммуникации между участниками проекта, Architecture-as-Code позволяет хранить описание архитектуры в формальном виде, версионировать его и применять в автоматизированных процессах.

Целью выпускной квалификационной работы является разработка и обоснование архитектуры автоматизированной платформы управления инфраструктурой с поддержкой подхода Architecture-as-Code, обеспечивающей формирование desired state, получение actual state, обнаружение drift и инициирование reconcile для приведения инфраструктуры к целевому состоянию.

Для достижения поставленной цели необходимо решить следующие задачи:
\begin{enumerate}
    \item проанализировать проблему управления инфраструктурой в распределённых системах и определить роль drift detection в эксплуатационном контуре;
    \item обосновать применение Architecture-as-Code как источника desired state для платформы управления инфраструктурой;
    \item сформировать функциональные и нефункциональные требования к платформе;
    \item спроектировать архитектуру платформы и определить ответственность сервисов;
    \item реализовать компоненты платформы;
    \item провести проверку качества разработанного решения.
\end{enumerate}

Объектом исследования являются процессы управления инфраструктурой распределённых программных систем. Предметом исследования являются архитектурные и программные решения, обеспечивающие автоматизированное сопоставление desired state и actual state, обнаружение drift и выполнение reconcile-действий.

Методы исследования включают изучение предметной области, сравнительный анализ существующих инструментов, проектирование программной архитектуры, прототипирование, разработку микросервисных компонентов, модульное тестирование и экспериментальную оценку поведения системы.

Практическая значимость работы состоит в снижении объёма ручных операций при контроле состояния инфраструктуры и в повышении наблюдаемости расхождений между архитектурным описанием и фактическим состоянием управляемых хостов. Разработанная платформа может использоваться как основа для дальнейшего развития внутреннего control plane, в котором архитектурная модель становится не только документацией, но и источником данных для автоматизированного управления.

Структура работы соответствует логике исследования и разработки. В первой главе рассматриваются теоретические основы управления инфраструктурой, понятия desired state, actual state, drift и reconcile, а также анализируются существующие инструменты и подходы. Во второй главе выполняется проектирование автоматизированной платформы управления инфраструктурой. В третьей главе описывается реализация основных компонентов платформы. В четвёртой главе приводится исследование и оценка качества разработанного решения. В пятой главе рассматриваются практическая значимость, ограничения текущего MVP и направления дальнейшего развития. В заключении формулируются основные результаты работы.
